\begin{abstract}
K-Waay的\batchreceive 构造规定，同一次调用中的不同输入应分别对应不同参与方。为考察这一条件在批处理接纳边界上的语义作用，本文对三个共享发送来源与顺序处理生命周期的固定两槽Tamarin符号模型进行受控比较，分别研究宽松接纳、消息级限制和参与方级限制。已有执行记录及后续独立复跑表明：在宽松接纳下，一个匹配的发送来源可支持同一批次上下文中的两个接收事件；消息级限制虽能保持消息区分和作用域内的来源单射性，同一参与方产生的两条不同消息仍可被共同接纳；参与方级限制直接约束所收集条目的参与方坐标，在模型接收事件边界上保持参与方区分，同时保留可达的有效异参与方批次。该比较说明，消息区分和发送实例区分不能替代批内参与方区分，批次组合边界需要显式维护参与方关系。分析以批次组合可控、新鲜符号消息和理想化精确来源匹配为前提，结论止于固定两槽模型中的\receiveraccept 事件，不构成对完整K-Waay执行的精化定理，也不表示密码原语被攻破或真实部署存在漏洞。后续独立复跑保留了证明器输出和导出轨迹，但并非对历史执行记录的恢复。
\end{abstract}
